% =============================================================================
% ANEXO S — Matriz de Cumplimiento ISO/IEC/IEEE 26515 y 15289
% =============================================================================
\chapter{Matriz de Cumplimiento Normativo}
\label{cap:anexo_iso}

Este anexo mapea las actividades exigidas por \textbf{ISO/IEC/IEEE 26515:2024} (desarrollo de
información en entornos ágiles) y \textbf{ISO/IEC/IEEE 15289:2024} (productos de información del
ciclo de vida) con las secciones de esta memoria que las evidencian.\par

\section{ISO/IEC/IEEE 26515 — Información en Entornos Ágiles}
\label{sec:anexo_iso_26515}
\renewcommand{\arraystretch}{1.3}
\fontsize{8.8}{10.5}\selectfont\sffamily
\noindent
\begin{tabularx}{\textwidth}{|X|X|c|}
\hline
\rowcolor{colorPrimario}
\textcolor{white}{\textbf{Actividad exigida}} & \textcolor{white}{\textbf{Evidencia en la memoria}} & \textcolor{white}{\textbf{Ref.}} \\ \hline
Planificación de la información por iteración & Sección de Gestión de la Información en cada sprint. & \ref{sec:s1_gestion} \\ \hline
\rowcolor{grisTecnico}
Información desarrollada en paralelo al código & \textit{Documentation as Code}; spec OpenAPI viva. & \ref{ssec:s0_estrategia_doc} \\ \hline
Historias de usuario para la información & Backlog de Épicas e Historias. & \ref{ssec:s0_backlog} \\ \hline
\rowcolor{grisTecnico}
Documentación continua de incrementos & Documentación de endpoints por sprint. & \ref{ssec:s2_gi_endpoints} \\ \hline
Revisión técnica de la información & \textit{Technical review} del Sprint 4. & \ref{ssec:s4_gi_review} \\ \hline
\rowcolor{grisTecnico}
Pruebas de usabilidad de la información & Validación de manuales y diagramas. & \ref{ssec:s4_gi_usabilidad} \\ \hline
Empaquetado de la información con el producto & Base de conocimiento embebida en el repo. & \ref{ssec:s5_gi_empaquetado} \\ \hline
\end{tabularx}\normalfont\normalsize
\par\vspace{0.3cm}

\section{ISO/IEC/IEEE 15289 — Productos de Información}
\label{sec:anexo_iso_15289}
\renewcommand{\arraystretch}{1.3}
\fontsize{8.8}{10.5}\selectfont\sffamily
\noindent
\begin{tabularx}{\textwidth}{|X|X|c|}
\hline
\rowcolor{colorPrimario}
\textcolor{white}{\textbf{Producto de información}} & \textcolor{white}{\textbf{Ubicación}} & \textcolor{white}{\textbf{Ref.}} \\ \hline
Descripción de la arquitectura & Modelo C4 (contexto, contenedores, componentes). & \ref{sec:s0_arquitectura} \\ \hline
\rowcolor{grisTecnico}
Especificación de la base de datos & Diccionario de datos y migraciones. & \ref{cap:anexo_datos} \\ \hline
Especificación de interfaces (API) & Endpoints, payloads y validaciones. & \ref{cap:anexo_ep_detalle} \\ \hline
\rowcolor{grisTecnico}
Registros de verificación y validación & Casos de prueba E2E. & \ref{cap:anexo_e2e} \\ \hline
Historial de revisiones & Control de versiones del documento. & --- \\ \hline
\rowcolor{grisTecnico}
Matriz de trazabilidad & Épicas $\leftrightarrow$ sprints $\leftrightarrow$ capítulos. & \ref{cap:anexo_trazabilidad} \\ \hline
Notas de la versión & Release Notes 2.0.0. & \ref{ssec:s5_gi_releasenotes} \\ \hline
\end{tabularx}\normalfont\normalsize

\begin{NotaTecnica}
La cobertura normativa es completa: cada actividad de ISO 26515 y cada producto de ISO 15289
tiene al menos una evidencia trazable en esta memoria o en sus anexos, versionada junto al
código fuente conforme al principio de \textit{Documentation as Code}.
\end{NotaTecnica}

\clearpage
